\begin{center}
    {\LARGE \textbf{2012物理化学}} \\[2ex]
    {\large 时量180分钟 \quad 满分150分}
\end{center}

\vspace{0.5cm}
\noindent\textbf{注意事项:}
\begin{enumerate}[label=\arabic*., parsep=0pt, itemsep=0pt]
    \item 答题前,考生先将自己的姓名、学号、考场号、座位号填写在答题卡上,将条形码准确粘贴在考生信息条形码粘贴区。
    \item 考试结束后,将本试卷与答题纸一并收回。
    \item 回答各个问题时,将答案写在答题纸上,写在本试卷上无效。
\end{enumerate}

\vspace{0.5cm}
\noindent\textbf{一、选择题（30 pts.）}

\begin{enumerate}[label=\arabic*., itemsep=1.5ex]
    \item 在温度$T$时，纯液体$\ce{A}$的饱和蒸气压为$p_\mathrm{A}^*$，化学势为$\mu_\mathrm{A}^*$，并且已知在$p^\ominus$压力下的凝固点为$T_\mathrm{f}^*$，当$\ce{A}$中溶入少量与$\ce{A}$不形成固态溶液的溶质而形成稀溶液时，上述三物理量分别为$p_\mathrm{A}$、$\mu_\mathrm{A}$、$T_\mathrm{f}$，则
    \begin{tasks}(2)
        \task $p_\mathrm{A}^* < p_\mathrm{A}$，$\mu_\mathrm{A}^* < \mu_\mathrm{A}$，$T_\mathrm{f}^* < T_\mathrm{f}$
        \task $p_\mathrm{A}^* > p_\mathrm{A}$，$\mu_\mathrm{A}^* < \mu_\mathrm{A}$，$T_\mathrm{f}^* < T_\mathrm{f}$
        \task $p_\mathrm{A}^* < p_\mathrm{A}$，$\mu_\mathrm{A}^* < \mu_\mathrm{A}$，$T_\mathrm{f}^* > T_\mathrm{f}$
        \task $p_\mathrm{A}^* > p_\mathrm{A}$，$\mu_\mathrm{A}^* > \mu_\mathrm{A}$，$T_\mathrm{f}^* > T_\mathrm{f}$
    \end{tasks}
    \item $\ce{Ag2O}$分解可用下面两个计量方程之一表示，其相应的平衡常数也一并列出：
    \[
    \ce{Ag2O(s) -> 2Ag(s) + \frac{1}{2}O2(g)} \quad K_p(1)
    \]
    \[
    \ce{2Ag2O(s) -> 4Ag(s) + O2(g)} \quad K_p(2)
    \]
    设气相为理想气体，且已知反应是吸热的，试判断下列结论哪个是正确的？
    \begin{tasks}(2)
        \task $K_p^\ominus(2)=K_p^{\ominus\frac{1}{2}}(1)$
        \task $K_p^\ominus(2)$随温度的升高而增大
        \task $K_p^\ominus(2)=K_p^\ominus(1)$
        \task $\ce{O2}$平衡压力与计量方程写法无关
    \end{tasks}
    \item $25^\circ\mathrm{C}$时，$1\ \mathrm{mol}$理想气体等温膨胀，压力从$10p^\ominus$变到$p^\ominus$，体系吉布斯自由能变化多少？
    \begin{tasks}(4)
        \task $0.04\ \mathrm{kJ}$
        \task $-12.4\ \mathrm{kJ}$
        \task $1.24\ \mathrm{kJ}$
        \task $-5.70\ \mathrm{kJ}$
    \end{tasks}
    \item 下列电池的电动势，哪个与$\ce{Br^-}$的活度无关？
    \begin{tasks}(1)
        \task $\ce{Ag(s) | AgBr(s) | KBr(aq) | Br2(l), Pt}$
        \task $\ce{Zn(s) | ZnBr2(aq) | Br2(l), Pt}$
        \task $\ce{Pt, H2(g) | HBr(aq) | Br2(l), Pt}$
        \task $\ce{Hg(l) | Hg2Br2(s) | KBr(aq) \parallel AgNO3(aq) | Ag(s)}$
    \end{tasks}
    \item 体系的状态改变了，其内能值
    \begin{tasks}(4)
        \task 必定改变
        \task 必定不变
        \task 不一定变
        \task 状态与内能无关
    \end{tasks}
    \item 苯的正常沸点为$80^\circ\mathrm{C}$，估计它在沸点左右温度范围内，温度每改变$1^\circ\mathrm{C}$，蒸气压的变化百分率约为
    \begin{tasks}(4)
        \task $3\%$
        \task $13\%$
        \task $47\%$
        \task 难以确定
    \end{tasks}
    \item $298\ \mathrm{K}$时，反应为$\ce{Zn(s) + Fe^{2+}(aq) = Zn^{2+}(aq) + Fe(s)}$的电池的$E^\ominus$为$0.323\ \mathrm{V}$，则其平衡常数$K^\ominus$为
    \begin{tasks}(4)
        \task $2.89\times10^5$
        \task $8.46\times10^{10}$
        \task $5.53\times10^4$
        \task $2.35\times10^7$
    \end{tasks}
    \item 下述说法哪一种不正确？
    \begin{tasks}(1)
        \task 理想气体经绝热自由膨胀后，其内能变化为零
        \task 非理想气体经绝热自由膨胀后，其内能变化不一定为零
        \task 非理想气体经绝热膨胀后，其温度一定降低
        \task 非理想气体经一不可逆循环，其内能变化为零
    \end{tasks}
    \item 对于一个一级反应，如其半衰期$t_{1/2}$在$0.01\ \mathrm{s}$以下，即称为快速反应，此时它的速率常数$k$值在
    \begin{tasks}(4)
        \task $69.32\ \mathrm{s^{-1}}$以上
        \task $6.932\ \mathrm{s^{-1}}$以上
        \task $0.069\ 32\ \mathrm{s^{-1}}$以上
        \task $6.932\ \mathrm{s^{-1}}$以下
    \end{tasks}
    \item $2\ \mathrm{mol}$液态苯在其正常沸点($353.2\ \mathrm{K}$)和$101.325\ \mathrm{kPa}$下蒸发为苯蒸气，该过程的
    
    $\Delta_{\mathrm{vap}}A$等于
    \begin{tasks}(4)
        \task $23.48\ \mathrm{kJ}$
        \task $5.87\ \mathrm{kJ}$
        \task $2.94\ \mathrm{kJ}$
        \task $1.47\ \mathrm{kJ}$
    \end{tasks}
    \item 加入惰性气体对哪个反应能增大其平衡转化率？
    \begin{tasks}(1)
        \task $\ce{C6H5C2H5(g) = C6H5C2H3(g) + H2(g)}$
        \task $\ce{CO(g) + H2O(g) = CO2(g) + H2(g)}$
        \task $\ce{\frac{3}{2}H2(g) + \frac{1}{2}N2(g) = NH3(g)}$
        \task $\ce{CH3COOH(l) + C2H5OH(l) = H2O(l) + C2H5COOCH3(l)}$
    \end{tasks}
    \item 下述体系中的组分$\ce{B}$，选择假想标准态的是
    \begin{tasks}(1)
        \task 混合理想气体中的组分$\ce{B}$
        \task 混合非理想气体中的组分$\ce{B}$
        \task 理想溶液中的组分$\ce{B}$
        \task 稀溶液中的溶剂
    \end{tasks}
    \item 下列四个偏微商中，哪个不是化学势？
    \begin{tasks}(4)
        \task $(\partial U/\partial n_\mathrm{B})_{S,V,n_c}$
        \task $(\partial H/\partial n_\mathrm{B})_{S,p,n_c}$
        \task $(\partial F/\partial n_\mathrm{B})_{T,p,n_c}$
        \task $(\partial G/\partial n_\mathrm{B})_{T,p,n_c}$
    \end{tasks}
    \item 对于同一电解质的水溶液，当其浓度逐渐增加时，何种性质将随之增加？
    \begin{tasks}(2)
        \task 在稀溶液范围内的电导率
        \task 摩尔电导率
        \task 电解质的离子平均活度系数
        \task 离子淌度
    \end{tasks}
    \item $273\ \mathrm{K}$，$10p^\ominus$下，液态水和固态水（即冰）的化学势分别为$\mu(\mathrm{l})$和$\mu(\mathrm{s})$，两者的关系为
    \begin{tasks}(4)
        \task $\mu(\mathrm{l}) > \mu(\mathrm{s})$
        \task $\mu(\mathrm{l}) = \mu(\mathrm{s})$
        \task $\mu(\mathrm{l}) < \mu(\mathrm{s})$
        \task 不确定
    \end{tasks}
\end{enumerate}

\vspace{0.5cm}
\noindent\textbf{二、计算题(20 pts.)}

在$573.15\ \mathrm{K}$时,将$1\ \mathrm{mol}\ \ce{Ne}$(可视为理想气体)从$1000\ \mathrm{kPa}$经绝热可逆膨胀到$100\ \mathrm{kPa}$。求$Q$,$W$,$\Delta U$,$\Delta H$,$\Delta S$,$\Delta S_\text{孤立}$和$\Delta G$。

已知在$573.15\ \mathrm{K}$、$1000\ \mathrm{kPa}$下$\ce{Ne}$的摩尔熵$S_\mathrm{m}=144.2\ \mathrm{J\cdot K^{-1}\cdot mol^{-1}}$。

\vspace{0.5cm}
\noindent\textbf{三、计算题(15 pts.)}

$2p^\ominus$的水蒸气可以维持一段时间,但这是一种亚稳平衡,称作过饱和态,它可自发地凝聚,过程是:
\[
\ce{H2O(g, 100^\circ C, 202.650 kPa) -> H2O(l, 100^\circ C, 202.650 kPa)}
\]
求水的$\Delta H_\mathrm{m}$,$\Delta S_\mathrm{m}$,$\Delta G_\mathrm{m}$。已知水的摩尔汽化焓$\Delta H_\mathrm{m}^\ominus$为$40.60\ \mathrm{kJ\cdot mol^{-1}}$,假设水蒸气为理想气体,液态水是不可压缩的。

\vspace{0.5cm}
\noindent\textbf{四、计算题(15 pts.)}

在$413.15\ \mathrm{K}$,总压为$101.325\ \mathrm{kPa}$时,$\ce{C6H5Cl}$和$\ce{C6H5Br}$两液体组成理想溶液混合物,该溶液的组成及液面上蒸气的组成分别为$x_\mathrm{A}=0.595\ 6$,$y_\mathrm{A}=0.736\ 2$。试求在$413.15\ \mathrm{K}$,纯$\ce{C6H5Cl}$和纯$\ce{C6H5Br}$的饱和蒸气压。

\vspace{0.5cm}
\noindent\textbf{五、计算题(20 pts.)}

$\ce{Mg}$(熔点$924\mathrm{K}$)和$\ce{Zn}$(熔点$692\mathrm{K}$)的相图具有两个低共熔点,一个为$641\ \mathrm{K}$($\ce{Mg}$的质量分数为$0.032$),另一个为$620\ \mathrm{K}$($\ce{Mg}$的质量分数为$0.49$),在体系的熔点曲线上有一个最高点$863\ \mathrm{K}$($\ce{Mg}$的质量分数为$0.157$)。
\begin{enumerate}[label=(\arabic*), itemsep=1ex]
    \item 绘出$\ce{Mg}$和$\ce{Zn}$的$T-x$图,并标明各区中的相；
    \item 分别指出含$\ce{Mg}$的质量分数为$0.80$和$0.30$的两个混合物从$973\ \mathrm{K}$冷却到$573\ \mathrm{K}$步冷过程中的相变,并根据相律予以说明；
    \item 绘出含$\ce{Mg}$的质量分数为$0.49$的熔化物的步冷曲线。
\end{enumerate}

\vspace{0.5cm}
\noindent\textbf{六、计算题(15 pts.)}

半导体工业为了获得氧含量不大于$1\times10^{-6}$的高纯氢,在$298\ \mathrm{K}$、$101\ \mathrm{kPa}$下让电解水制得的氢气($99.5\%\ \ce{H2}$,$0.5\%\ \ce{O2}$)通过催化剂,发生反应$\ce{2H2(g) + O2(g) = 2H2O(g)}$达到消除氧的目的。试问反应后氢气的纯度能否达到标准要求?已知$\Delta_\mathrm{f}G_\mathrm{m}^\ominus(\ce{H2O,g})=-228.59\ \mathrm{kJ\cdot mol^{-1}}$。

\vspace{0.5cm}
\noindent\textbf{七、计算题(15 pts.)}

在$298\ \mathrm{K}$时有电池:
$\ce{Pt,H2(p^\ominus) | HCl(0.01mol\cdot Kg^{-1}) | AgCl(s),Ag(s)}$
已知$\ce{AgCl}$的$K_{\mathrm{sp}}=1.745\times10^{-10}$。
\begin{enumerate}[label=(\arabic*), itemsep=1ex]
    \item 写出电池反应；
    \item 计算$E^\ominus(\ce{AgCl/Ag})$；
    \item 计算电动势$E$。
\end{enumerate}

\vspace{0.5cm}
\noindent\textbf{八、计算题(20 pts.)}

$25^\circ\mathrm{C}$时,电池$\ce{Ag | AgCl(s) | HCl(0.1\ mol\cdot kg^{-1}) | Cl2(p^\ominus),Pt}$的电池电动势$E$为$1.107\ \mathrm{V}$。

已知：$\ce{AgCl}$在$25^\circ\mathrm{C}$时的标准生成焓和标准生成吉布斯自由能分别为$-127.03\ \mathrm{kJ\cdot mol^{-1}}$和

$-106.83\ \mathrm{kJ\cdot mol^{-1}}$,试计算$25^\circ\mathrm{C}$时:
\begin{enumerate}[label=(\arabic*), itemsep=1ex]
    \item 该电池反应的$\Delta_\mathrm{r}S_\mathrm{m}^\ominus$；
    \item 该电池的可逆热效应；
    \item 电池的温度系数。
\end{enumerate}